% Shared preamble of the RMC kernel write-ups (% Shared preamble of the RMC kernel write-ups (% Shared preamble of the RMC kernel write-ups (% Shared preamble of the RMC kernel write-ups (\input{preamble} after \documentclass).
\usepackage[margin=1in]{geometry}
\usepackage{amsmath,amssymb,bm}
\usepackage{graphicx}
\usepackage{booktabs}
\usepackage{hyperref}

\newcommand{\dd}{\,\mathrm{d}}
\newcommand{\Ein}{E}
\newcommand{\Eout}{E'}
\newcommand{\Ecm}{E_{c}}
\newcommand{\mucm}{\mu_{c}}
\newcommand{\mulab}{\mu_0}

% Common closing section: how the verification figures are produced
\newcommand{\verificationmethod}{%
Each figure is produced by \texttt{verification/verify\_laws.py}. For an incident
energy $\Ein$ along $+z$, $10^6$ collisions are sampled with MC/DC's own reaction
sampler (\texttt{sample\_elastic\_scattering}, \texttt{sample\_inelastic\_scattering}
or \texttt{sample\_fission}), and every emitted neutron is histogrammed in
$48 \times 24$ bins of $(\Eout, \mulab)$, with $\mulab$ the lab cosine to $+z$. The
inverted kernel used by RMC is integrated over the same bins (Gauss--Legendre panels;
two-body lines are split where $\mulab^*(\Eout)$ crosses a cosine edge), multiplied by
the number of collisions, and compared with a $\chi^2$ test over bins with at least 5
expected counts (the rest pooled into one bin). The tables list $\chi^2/\mathrm{dof}$,
its $p$-value, the emitted neutrons per collision (sampled vs.\ the integral of the
inverted kernel), and the largest relative change of a bin integral when the
quadrature panels are doubled. Left panels: outgoing-energy marginal; middle:
lab-cosine marginal; right: per-bin $z = (\text{sampled} - \text{inverted}) /
\sqrt{\text{inverted}}$.}
 after \documentclass).
\usepackage[margin=1in]{geometry}
\usepackage{amsmath,amssymb,bm}
\usepackage{graphicx}
\usepackage{booktabs}
\usepackage{hyperref}

\newcommand{\dd}{\,\mathrm{d}}
\newcommand{\Ein}{E}
\newcommand{\Eout}{E'}
\newcommand{\Ecm}{E_{c}}
\newcommand{\mucm}{\mu_{c}}
\newcommand{\mulab}{\mu_0}

% Common closing section: how the verification figures are produced
\newcommand{\verificationmethod}{%
Each figure is produced by \texttt{verification/verify\_laws.py}. For an incident
energy $\Ein$ along $+z$, $10^6$ collisions are sampled with MC/DC's own reaction
sampler (\texttt{sample\_elastic\_scattering}, \texttt{sample\_inelastic\_scattering}
or \texttt{sample\_fission}), and every emitted neutron is histogrammed in
$48 \times 24$ bins of $(\Eout, \mulab)$, with $\mulab$ the lab cosine to $+z$. The
inverted kernel used by RMC is integrated over the same bins (Gauss--Legendre panels;
two-body lines are split where $\mulab^*(\Eout)$ crosses a cosine edge), multiplied by
the number of collisions, and compared with a $\chi^2$ test over bins with at least 5
expected counts (the rest pooled into one bin). The tables list $\chi^2/\mathrm{dof}$,
its $p$-value, the emitted neutrons per collision (sampled vs.\ the integral of the
inverted kernel), and the largest relative change of a bin integral when the
quadrature panels are doubled. Left panels: outgoing-energy marginal; middle:
lab-cosine marginal; right: per-bin $z = (\text{sampled} - \text{inverted}) /
\sqrt{\text{inverted}}$.}
 after \documentclass).
\usepackage[margin=1in]{geometry}
\usepackage{amsmath,amssymb,bm}
\usepackage{graphicx}
\usepackage{booktabs}
\usepackage{hyperref}

\newcommand{\dd}{\,\mathrm{d}}
\newcommand{\Ein}{E}
\newcommand{\Eout}{E'}
\newcommand{\Ecm}{E_{c}}
\newcommand{\mucm}{\mu_{c}}
\newcommand{\mulab}{\mu_0}

% Common closing section: how the verification figures are produced
\newcommand{\verificationmethod}{%
Each figure is produced by \texttt{verification/verify\_laws.py}. For an incident
energy $\Ein$ along $+z$, $10^6$ collisions are sampled with MC/DC's own reaction
sampler (\texttt{sample\_elastic\_scattering}, \texttt{sample\_inelastic\_scattering}
or \texttt{sample\_fission}), and every emitted neutron is histogrammed in
$48 \times 24$ bins of $(\Eout, \mulab)$, with $\mulab$ the lab cosine to $+z$. The
inverted kernel used by RMC is integrated over the same bins (Gauss--Legendre panels;
two-body lines are split where $\mulab^*(\Eout)$ crosses a cosine edge), multiplied by
the number of collisions, and compared with a $\chi^2$ test over bins with at least 5
expected counts (the rest pooled into one bin). The tables list $\chi^2/\mathrm{dof}$,
its $p$-value, the emitted neutrons per collision (sampled vs.\ the integral of the
inverted kernel), and the largest relative change of a bin integral when the
quadrature panels are doubled. Left panels: outgoing-energy marginal; middle:
lab-cosine marginal; right: per-bin $z = (\text{sampled} - \text{inverted}) /
\sqrt{\text{inverted}}$.}
 after \documentclass).
\usepackage[margin=1in]{geometry}
\usepackage{amsmath,amssymb,bm}
\usepackage{graphicx}
\usepackage{booktabs}
\usepackage{hyperref}

\newcommand{\dd}{\,\mathrm{d}}
\newcommand{\Ein}{E}
\newcommand{\Eout}{E'}
\newcommand{\Ecm}{E_{c}}
\newcommand{\mucm}{\mu_{c}}
\newcommand{\mulab}{\mu_0}

% Common closing section: how the verification figures are produced
\newcommand{\verificationmethod}{%
Each figure is produced by \texttt{verification/verify\_laws.py}. For an incident
energy $\Ein$ along $+z$, $10^6$ collisions are sampled with MC/DC's own reaction
sampler (\texttt{sample\_elastic\_scattering}, \texttt{sample\_inelastic\_scattering}
or \texttt{sample\_fission}), and every emitted neutron is histogrammed in
$48 \times 24$ bins of $(\Eout, \mulab)$, with $\mulab$ the lab cosine to $+z$. The
inverted kernel used by RMC is integrated over the same bins (Gauss--Legendre panels;
two-body lines are split where $\mulab^*(\Eout)$ crosses a cosine edge), multiplied by
the number of collisions, and compared with a $\chi^2$ test over bins with at least 5
expected counts (the rest pooled into one bin). The tables list $\chi^2/\mathrm{dof}$,
its $p$-value, the emitted neutrons per collision (sampled vs.\ the integral of the
inverted kernel), and the largest relative change of a bin integral when the
quadrature panels are doubled. Left panels: outgoing-energy marginal; middle:
lab-cosine marginal; right: per-bin $z = (\text{sampled} - \text{inverted}) /
\sqrt{\text{inverted}}$.}
